
\subsubsection{Spurious Correlations and Group Robustness}

The study of spurious correlations in deep learning traces to observations that models exploit dataset biases rather than causal features [Geirhos et al., 2020; Zhang et al., 2017]. Sagawa et al. [2020] formalize this via group robustness: in Waterbirds and CelebA, models achieving high average accuracy fail badly on minority groups whose spurious attribute contradicts the majority correlation. Group DRO minimizes worst-group loss but requires group annotations at training time. Our work is complementary: we characterize how much spurious content enters frozen representations \textit{before} any downstream training, as a function of pretraining paradigm.

Kirichenko et al. [2022] demonstrate that ERM features are sufficient for robustness after group-balanced last-layer retraining (Deep Feature Reweighting, DFR), motivating direct study of frozen representation quality. However, DFR focuses entirely on supervised ERM and does not compare to SSL paradigms --- leaving open whether ERM's spurious feature content is typical or atypical among pretraining objectives.

Izmailov et al. [2022] provide a study of feature learning under spurious correlations that includes comparisons involving DINO and Waterbirds. Their analysis is closest to ours in scope, but does not use a controlled spurious/task ratio metric and does not perform a four-paradigm comparison with Bonferroni-corrected pairwise tests [citation requires verification before submission].

\subsubsection{Self-Supervised and Contrastive Representation Learning}

Contrastive SSL methods [Chen et al., 2020 (SimCLR); He et al., 2020 (MoCo)] learn representations by maximizing agreement between augmented views of the same image. DINO \citep{Caron2021DINO} uses self-distillation with a momentum teacher, producing emergent class-discriminative attention maps. BarlowTwins \citep{Zbontar2021BarlowTwins} optimizes redundancy reduction across cross-correlation matrices. All four paradigms use ResNet-50 as a standard evaluation backbone.

These methods are primarily evaluated on ImageNet linear probe accuracy [Chen et al., 2020; Caron et al., 2021], which measures task-relevant encoding but says nothing about spurious feature content. Robinson et al. [2021] show that contrastive learning can encode shortcut features, particularly when spurious attributes are stable across augmentations [citation requires verification]. Wen et al. [2021] analyze what features contrastive objectives encode as a function of augmentation design [citation requires verification]. These works establish that SSL is not spurious-feature-free, but do not compare to supervised baselines under controlled conditions.

\subsubsection{Distribution Shift Benchmarks}

WILDS \citep{Koh2021WILDS} and DomainBed \citep{Gulrajani2021DomainBed} evaluate fine-tuned models, making it impossible to attribute performance differences to pretraining representation quality versus fine-tuning procedure. We isolate the pretraining effect by freezing the backbone throughout.

\subsubsection{Augmentation and Feature Learning}

Shah et al. [2020] characterize SGD's ''simplicity bias'' --- the tendency to prefer low-complexity predictors --- connecting to our finding that ERM encodes spurious features aggressively when they are predictive under the training distribution. Zimmermann et al. [2021] show theoretically that contrastive learning approximately inverts the data generating process. Our empirical finding --- that label-agnostic contrastive SSL encodes spurious features less than label-supervised ERM --- is consistent with this framework when spurious attributes are not part of the true latent structure.

